%% AgenticDev 2026 -- workshop summary for ACM SIGSOFT Software Engineering Notes (SEN)
%%
%% TEMPORARY PREAMBLE. SEN requires its own class, acm_sen_article.cls
%% (https://www2.sigsoft.org/sen/guideline/). Until the official class and
%% example are added to the repository, this file builds with a neutral
%% two-column article layout so the content can be drafted. Only the
%% preamble and \maketitle block below will change when we switch.
\documentclass[10pt,twocolumn]{article}

\usepackage[T1]{fontenc}
\usepackage[utf8]{inputenc}
\usepackage[british]{babel}
\usepackage{mathptmx}
\usepackage[margin=0.75in]{geometry}
\usepackage{graphicx}
\usepackage{booktabs}
\usepackage{xcolor}
\usepackage[numbers,sort&compress]{natbib}
\usepackage[hidelinks]{hyperref}

% Shared macros (keep independent of the document class).
\newcommand{\agdev}{AgenticDev~2026}
% Visible drafting notes; set \drafttrue/\draftfalse to show/hide.
\newif\ifdraft \drafttrue
\newcommand{\todo}[1]{\ifdraft{\color{red}\textbf{[TODO:} #1\textbf{]}}\fi}
\newcommand{\note}[2]{\ifdraft{\color{blue}\textbf{[#1:} #2\textbf{]}}\fi}


\title{Agentic AI for Next-Generation Software Development:\\
  Summary of the First AgenticDev Workshop (AgenticDev 2026)}

\author{%
  Andrea Rosani\\ Free University of Bozen-Bolzano, Italy\\ \texttt{andrea.rosani@unibz.it}
  \and
  Atta Badii\\ University of Reading, UK\\ \texttt{sis13i@reading.ac.uk}
  \and
  Ali Atalay\\ AI4SEC O\"U, Estonia\\ \texttt{ali.atalay@ai4sec.eu}
  \and
  Simos Gerasimou\\ Cyprus University of Technology, Cyprus\\ \texttt{simos.gerasimou@cut.ac.cy}
  \and
  Jean-Marie Mottu\\ Universit\'e de Nantes -- LS2N -- IMT Atlantique, France\\ \texttt{jean-marie.mottu@univ-nantes.fr}
  \and
  Giuseppe Di Fatta\\ Free University of Bozen-Bolzano, Italy\\ \texttt{giuseppe.difatta@unibz.it}
}
\date{}

\begin{document}
\maketitle

\begin{abstract}
% SEN: the abstract appears in the printed issue; the full summary in the ACM DL.
% Structure it to reflect the key points: scope, programme, research directions, outlook.
\todo{Abstract (approx.\ 150--200 words): the First International Workshop on Agentic AI for Next-Generation Software Development (\agdev), co-located with ASE~2026 in Munich on 12~October~2026; aims; submissions and accepted contributions; research directions identified; outlook.}

\end{abstract}

\noindent\textbf{Keywords:} agentic AI, large language models, AI-assisted software engineering, multi-agent systems, coding agents, human--AI collaboration
% TODO: add ACM CCS concepts (https://dl.acm.org/ccs), as SEN guidelines recommend.

\section{Introduction}
\label{sec:introduction}

\todo{Motivation: from LLM-based code completion to agentic systems that plan, use tools and collaborate across the software lifecycle.}

\todo{Aim and scope of \agdev: agent architectures and coordination, human--AI collaboration, integration into engineering workflows; trustworthiness, oversight, verification and validation, explainability, empirical evaluation.}

\todo{Link to the four Horizon Europe projects that inspired the workshop: AI4SWEng, ASTIR, HIVEMIND, MOSAICO.}

\todo{Contribution and structure of this summary.}

\section{The Workshop}
\label{sec:workshop}

\subsection{Organisation and Review Process}
\todo{Date and venue; organising committee; programme committee; submission categories (long, short, demo); 84 submissions; at least two reviews each; accepted contributions (11 long, 8 short, extended abstracts) -- reconcile the 25 stated in the foreword with the 26 papers in the proceedings.}

\subsection{Programme}
\todo{Keynote(s), paper sessions, the AGenSE Alliance panel, attendance.}

\section{Agentic AI for Software Engineering: State of the Art}
\label{sec:sota}

\todo{Concise review positioning the workshop. Candidate threads, to be agreed:
(i)~from code LLMs to autonomous coding agents (agent--computer interfaces, tool use, planning);
(ii)~repository-level benchmarks and evaluation methodology;
(iii)~multi-agent architectures and orchestration;
(iv)~agents across the lifecycle: requirements, design and architecture, testing, maintenance;
(v)~trust, safety, security and human oversight;
(vi)~cost and efficiency of agentic workflows.}

\section{Research Directions from the Workshop}
\label{sec:directions}

\todo{Group the accepted contributions by the research directions they address. For each direction: the problem, how the workshop papers tackle it, and the key findings. Paper list and abstracts: \texttt{notes/accepted-papers.md}; citation keys \texttt{agdev26:<first-author>}.}

% \subsection{Direction 1}
% \subsection{Direction 2}
% ...

\section{Panel: The AGenSE Alliance}
\label{sec:panel}

\todo{Summary of the panel ``The AGenSE Alliance: Shaping the Future of Agentic Software Engineering'': panellists, questions posed, main positions and points of (dis)agreement.}

\section{Open Challenges and Future Directions}
\label{sec:future}

\todo{Future developments the papers and the discussion point to as most promising for the field; open problems; a research agenda.}

\section{Conclusions}
\label{sec:conclusion}

\todo{Takeaways and plans for future editions of AgenticDev.}


\section*{Acknowledgements}
\todo{Thank authors, programme committee, panellists and the ASE~2026 organisers. Acknowledge Horizon Europe funding (AI4SWEng, ASTIR, HIVEMIND, MOSAICO -- grant agreement numbers).}


\bibliographystyle{abbrvnat}
\bibliography{references}

\end{document}
